%%%%%%%%%%%%%%%%%%%%%%%%%%%%%%%%%%%%%%%%%%%%%%%%%%%%%%%%%%%%%%%%%%%%%%%%%%%%%
%  MLSys 2027 submission -- 2 column, <=10 pages excl. references
%  Compile: pdflatex paper && bibtex paper && pdflatex paper && pdflatex paper
%%%%%%%%%%%%%%%%%%%%%%%%%%%%%%%%%%%%%%%%%%%%%%%%%%%%%%%%%%%%%%%%%%%%%%%%%%%%%
\documentclass{article}
\usepackage{mlsys2025}                 % anonymous by default; \usepackage[accepted]{mlsys2025} for camera-ready

\usepackage{amsmath,amssymb,amsfonts,amsthm}
\usepackage{graphicx}
\usepackage{booktabs}
\usepackage{multirow}
\usepackage{array}
\usepackage{xcolor}
\usepackage{colortbl}
\usepackage{tikz}
\usepackage{pgfplots}
\pgfplotsset{compat=1.17}
\usetikzlibrary{arrows.meta,positioning}
\usepackage{algorithm}
\usepackage{algorithmic}
\usepackage{hyperref}
\usepackage{url}

\definecolor{hl}{RGB}{232,238,250}
\definecolor{bad}{RGB}{168,62,42}
\definecolor{keep}{RGB}{38,70,142}

\newtheorem{proposition}{Proposition}
\newtheorem{remark}{Remark}

\mlsystitlerunning{Learned Decompression for Capacity-Bounded Agent Memory}

\begin{document}

\twocolumn[
\mlsystitle{Cue-Conditioned Reconstruction: Learned Decompression for
Capacity-Bounded Agent Memory}

\begin{mlsysauthorlist}
\mlsysauthor{Anonymous}{anon}
\end{mlsysauthorlist}
\mlsysaffiliation{anon}{Anonymous Institution}
\mlsyscorrespondingauthor{Anonymous}{anon@example.com}
\mlsyskeywords{agent memory, lossy compression, side information, calibration}

\vskip 0.3in

\begin{abstract}
An agent that must hold its memory within a fixed budget has to destroy
information, and the question is which representation of the survivors best
supports later recall. Current systems compress heuristically -- truncate,
summarise, or delete -- with no model of what a query will need. We formulate
memory compaction as \emph{lossy source coding with side information at the
decoder}: the encoder compacts a trace without knowing the future query, while
the decoder sees contextual cues the encoder did not. This is the Wyner-Ziv
setting, and it implies a compaction rule that differs sharply from
summarisation: discard preferentially what the cues will make recoverable, and
spend bits on what they will not.

We instantiate this as a trained encoder-decoder pair over an agent memory
store, with a decoder that emits a calibrated confidence and abstains rather
than reconstruct when the cues are insufficient. Training requires no
annotation: every compaction produces a (compacted trace, cues, original)
triple automatically, so the compressor generates its own supervision.

We first measure the ceiling. On a 1{,}115-memory corpus compacted to gist, the
answer to a question survives compaction for $46\%$ of questions, is absent but
latent in neighbouring memories for $7\%$, and is unrecoverable for $48\%$. A
learned decoder can therefore address at most the middle band, and any system
that produces fluent output on the remaining $48\%$ is confabulating by
construction. We argue this decomposition, not reconstruction accuracy, is the
correct way to evaluate learned memory decompression, and we report abstention
quality as the primary metric.
\end{abstract}
]
\printAffiliationsAndNotice{}

% =====================================================================
\section{Introduction}
\label{sec:intro}

% MOTIVATION
% - persistent agents exceed any fixed budget; capacity is the binding constraint
% - current compaction is heuristic: truncate / summarise / delete
% - none of it models what a future query will need
% - the ICLR companion work shows demotion beats deletion, but uses a
%   hand-written tier ladder; this paper asks what the ladder should be

\subsection{Contributions}
% (i)   Wyner-Ziv formulation of memory compaction
% (ii)  self-supervised training scheme with zero annotation cost
% (iii) the 46/7/48 recoverability decomposition as an evaluation protocol
% (iv)  a calibrated decoder whose headline metric is abstention, not fluency

% =====================================================================
\section{Memory compaction as Wyner-Ziv coding}
\label{sec:theory}

% Setup: trace X, compacted \hat{X}_\ell, cues Y available only at decode time.
% Encoder does not see Y. Decoder sees (\hat{X}_\ell, Y).
% R_{WZ}(D) <= R_{X|Y}(D): side information at the decoder alone nearly
% matches having it at both ends.
%
% The design consequence, which is the point of this section:
%   argmin over what-to-drop is NOT "least important" but "most predictable
%   from Y". A summariser optimises the wrong objective.
%
% Worked example from the corpus:
%   full: "Caroline attended an LGBTQ support group on 7 May 2023 and found
%          the transgender stories inspiring"
%   a summariser drops the date (long, looks incidental)
%   a Wyner-Ziv encoder keeps the date (unpredictable from cues) and drops
%   "found the stories inspiring" (predictable from topic + tag)

% =====================================================================
\section{The recoverability ceiling}
\label{sec:ceiling}

% This section must come BEFORE the method. It bounds what any decoder can do
% and it is the paper's most defensible contribution.

\begin{table}[t]
\caption{Where the answer goes when a memory is compacted to gist. Measured on
147 evidence-annotated questions over a 1{,}115-memory store.}
\label{tab:ceiling}
\centering\small
\begin{tabular}{@{}lrl@{}}
\toprule
\textbf{After compaction} & \textbf{Share} & \textbf{Implication} \\
\midrule
answer survives in the gist        & 46\% & decoder not needed \\
\rowcolor{hl}
lost, but latent in neighbours     & \phantom{0}7\% & \textbf{addressable} \\
unrecoverable                      & 48\% & \textcolor{bad}{must abstain} \\
\bottomrule
\end{tabular}
\end{table}

% Method for obtaining the table (must be exact and reproducible):
%   - gold answer tokens vs tokens present in gist / neighbour gists
%   - threshold 0.5 token overlap
%   - report sensitivity to that threshold in the appendix
%
% Consequence stated plainly: a decoder trained to always produce output will
% confabulate on 48% of queries. Abstention is not a safety add-on, it is the
% primary capability.

% =====================================================================
\section{Method}
\label{sec:method}

\subsection{Encoder: what to discard}
% E_\ell : X -> \hat{X}_\ell, rate-constrained, trained against decoder loss
% rather than a summarisation objective.

\subsection{Decoder: reconstruct or abstain}
% D : (\hat{X}_\ell, Y) -> (\tilde{X}, c)
% small instruction-tuned base (1-3B); cues = extracted cue set + k neighbour
% gists + encoding timestamp + topic.

\subsection{Escalation}
% retrieve at tier 0 -> widen to tier \ell -> decode -> abstain if c < tau.
% This is a cascade; the abstention threshold is the deployment knob.

\subsection{Objective}
\begin{equation}
\mathcal{L} = \mathcal{L}_{\text{recon}}
            + \lambda_1 \mathcal{L}_{\text{cal}}
            + \lambda_2 \mathcal{L}_{\text{abstain}}
            + \lambda_3 \mathcal{L}_{\text{rate}}
\label{eq:loss}
\end{equation}
% L_abstain is asymmetric: a confident wrong reconstruction costs >> an
% abstention. Report the ratio as a hyperparameter and sweep it.

% =====================================================================
\section{Training data}
\label{sec:data}

% The compressor generates its own supervision:
%   for each trace X and tier \ell:  (\hat{X}_\ell, Y, X)
% no annotation, scales with any conversational corpus.
%
% Cue sampling is the design choice that matters: VARY cue strength during
% training so the model learns the map from cue sufficiency to confidence,
% rather than assuming rich context.
%
% Corpora: LoCoMo, LongMemEval, MSC, plus non-dialogue sources
% (commit histories, issue threads) to test transfer.

% =====================================================================
\section{Evaluation protocol}
\label{sec:eval}

% Deliberately NOT reconstruction BLEU -- it rewards fluent confabulation.
%   1. faithful recovery rate on the addressable 7% band
%   2. confabulation rate on the unrecoverable 48%  <- primary safety metric
%   3. calibration (ECE) between c and correctness
%   4. end-to-end recall under a capacity bound, vs the heuristic tier ladder
%   5. systems metrics: decode latency, tokens/query, memory footprint
%      (MLSys audience cares about these as first-class results)

% =====================================================================
\section{Related work}
\label{sec:related}

% - agent memory systems and their compaction policies
% - Wyner-Ziv / distributed source coding
% - learned compression (neural codecs), semantic compression
% - calibration and selective prediction / abstention
% - retrieval-augmented generation and its hallucination literature

% =====================================================================
\section{Conclusion}
\label{sec:conclusion}

\bibliography{paper}
\bibliographystyle{mlsys2025}

\end{document}
